\documentclass[12pt]{article}
\usepackage[utf8]{inputenc}
\usepackage[russian]{babel}
\usepackage{amsthm,amsmath,amssymb}
\usepackage{fullpage}

 
\newtheorem{theorem}{Теорема}
\newtheorem{lemma}{Лемма}

\begin{document}

  \title{Продвинутые алгоритмы: алгоритмы для NP-трудных задач}
  \author{Кожагулов Рамазан}
  \date{}
  \maketitle


  \section*{Задача 0.1. Вершинное покрытие минимальной мощности — NP‑трудность}

  \paragraph{Условие.}
  Покажите, что задача о вершинном покрытии наименьшей мощности является NP‑трудной. Для этого достаточно доказать NP‑трудность соответствующей задачи разрешимости: языка VERTEXCOVER $= \{\langle G,k\rangle : G=(V,E)$ — неориентированный граф, и в $G$ есть вершинное покрытие размера не более $k\}$.

  \paragraph{Решение.}
  Сведём задачу о клике (язык CLIQUE $= \{\langle G,k\rangle : G$ содержит клику размера не менее $k\}$) к VERTEXCOVER за полиномиальное время. Пусть дан вход $\langle G=(V,E),k\rangle$. Построим дополнение графа $\overline{G}=(V,\overline{E})$, где $\overline{E}=\{\{u,v\} : u\neq v,\ \{u,v\}\notin E\}$, и положим $k' := |V|-k$. Определим функцию сведения $f(\langle G,k\rangle) := \langle \overline{G}, k'\rangle$.

  Корректность сведения основана на классическом эквиваленте: подмножество $S\subseteq V$ является кликой размера $k$ в $G$ тогда и только тогда, когда дополнение $V\setminus S$ является вершинным покрытием мощности $|V|-k$ в $\overline{G}$. Действительно, если $S$ — клика в $G$, то никакое ребро $\overline{G}$ не соединяет две вершины из $S$, значит, каждое ребро $\overline{G}$ инцидентно хотя бы одной вершине из $V\setminus S$, то есть $V\setminus S$ — вершинное покрытие. Обратно, если $X\subseteq V$ — вершинное покрытие в $\overline{G}$ мощности $|V|-k$, то на множестве $S:=V\setminus X$ нет рёбер $\overline{G}$, следовательно, $S$ образует клику размера $k$ в $G$.

  Размер и время построения полиномиальны: дополнение графа считается за $O(|V|^2)$, а вычисление $k'$ — константно. Поскольку $\textbf{CLIQUE}$ NP‑полон (NP‑complete), из сведения $\textbf{CLIQUE} \leq_p \textbf{VERTEXCOVER}$ следует, что $\textbf{VERTEXCOVER}$ NP‑труден. Отсюда оптимизационная версия задачи о вершинном покрытии также является NP‑трудной.

  \section*{Задача 0.2. Ациклический подграф — 1/2‑приближённый алгоритм}

  \paragraph{Условие.}
  Дано: ориентированный граф $G=(V,E)$. Требуется найти множество дуг $E^*\subseteq E$ максимальной мощности, чтобы подграф $(V,E^*)$ не содержал циклов.

  \paragraph{Алгоритм (приближение 1/2).}
  При случайном порядке вершин, для каждой дуги $e\in E$ независимо подбрасываем монету, с вероятностью $1/2$ помещаем $e$ в первое множество, иначе — во второе. Получаем два подграфа на тех же вершинах. Из них выбираем тот, который не содержит циклов и имеет больше дуг (если оба ацикличны — берём больший по числу дуг, если ацикличен только один — берём его). Математическое ожидание числа дуг в каждом из подграфов равно половине от числа дуг оптимального ациклического подграфа, поэтому такой выбор даёт коэффициент приближения $1/2$ по ожиданию.

  Детминизованный вариант той же гарантии. Упорядочим вершины произвольно. Рассмотрим две ориентации относительно этого порядка: все дуги, идущие слева направо (по возрастанию номеров), образуют ациклический подграф, все дуги, идущие справа налево, также образуют ациклический подграф. Выберем из этих двух подграфов тот, где дуг больше. Как минимум одна из этих двух сторон покрывает не менее половины дуг оптимального решения, что даёт ту же гарантию $1/2$ без случайности.

  \paragraph{Корректность.}
  В оптимальном решении дуги можно разбить на те, что согласованы с некоторым порядком обхода вершин, и обратные к нему. Хотя сам оптимальный порядок неизвестен, любой фиксированный порядок разбивает дуги пополам, не менее половины дуг направлено по порядку, не более половины — против. Подграф, состоящий из дуг, направленных по фиксированному порядку, ацикличен и содержит не менее половины дуг оптимального решения, аналогично для противоположного направления. Поэтому выбор лучшего из двух даёт коэффициент $1/2$.
  
  \section*{Задача 0.3. Минимальное максимальное паросочетание — 2‑приближённый алгоритм}

  \paragraph{Условие.}
  Дано: неориентированный граф $G=(V,E)$. Требуется найти паросочетание $E'\subseteq E$, которое является максимальным по включению (к $E'$ нельзя добавить ни одного ребра, сохранив паросочетание) и при этом имеет как можно меньше рёбер.

  \paragraph{Алгоритм (приближение 2).}
  Построим любое максимальное по включению паросочетание жадно: идём по рёбрам в любом порядке и добавляем ребро $(u,v)$ в ответ, если ни $u$, ни $v$ ещё не покрыты ранее добавленными рёбрами. В конце получаем паросочетание, к которому нельзя добавить новое ребро, то есть оно максимальное по включению.

  \paragraph{Корректность.}
  Обозначим через $\nu$ размер максимального (по мощности) паросочетания, а через $\text{OPT}$ — минимальный возможный размер среди всех максимальных по включению паросочетаний. Из известного факта следует, что любое максимальное по включению паросочетание имеет размер не меньше $\nu/2$, значит $\text{OPT}\ge \nu/2$. С другой стороны, полученное жадным алгоритмом паросочетание — это просто некоторое паросочетание, поэтому его размер не превосходит $\nu$. Следовательно, для результата $M$ жадного алгоритма верно $|M|\le \nu\le 2\,\text{OPT}$, то есть отношение не хуже $2$. Время работы — $O(|E|)$ при хранении пометки «покрыта/не покрыта» для каждой вершины.


  \section*{Задача 0.4. Вершинное покрытие через внутренние вершины поиска в глубину — 2‑приближённый алгоритм}

  \paragraph{Условие.}
  Дано: неориентированный граф $G=(V,E)$. Построить в $G$ дерево поиска в глубину и взять множество внутренних вершин $S$ этого дерева (вершины с хотя бы одним ребёнком). Доказать, что $S$ — вершинное покрытие и что $|S|\le 2\cdot \mathrm{OPT}$, где $\mathrm{OPT}$ — размер минимального вершинного покрытия.

  \paragraph{Алгоритм.}
  Выполним поиск в глубину, получим дерево и вернём в ответ множество всех внутренних вершин $S$ этого дерева.

  \paragraph{Корректность.}
  Во‑первых, $S$ — вершинное покрытие. Любое ребро либо является рёбром дерева (между родителем и ребёнком), тогда родитель внутренний и входит в $S$, либо это нетривиальное ребро к предку в дереве, а предок имеет хотя бы одного ребёнка, значит он тоже внутренний и входит в $S$. Следовательно, каждое ребро инцидентно вершине из $S$.

  Во‑вторых, $|S|\le 2\cdot \mathrm{OPT}$. Построим в дереве паросочетание $M$ размера не меньше $|S|/2$: обойдём вершины в порядке убывания глубины и, если у внутренней вершины есть ребёнок, ещё не инцидентный рёбрам $M$, добавим рёбра родитель–ребёнок в $M$. Это паросочетание, потому что каждый ребёнок выбирается не более одного раза, а родитель после добавления больше не используется. Каждая внутренняя вершина либо даёт одно ребро $M$, либо «прикрепляется» к ребру своего родителя, поэтому $|S|\le 2|M|$. Любое вершинное покрытие имеет размер не меньше размера любого паросочетания, значит $\mathrm{OPT}\ge |M|\ge |S|/2$, откуда $|S|\le 2\cdot \mathrm{OPT}$.

  \section*{Задача 0.5. Жадное покрытие вершинами по максимальной степени — гарантия и контрпример}

  \paragraph{Условие.}
  Дан неориентированный граф $G=(V,E)$. Рассмотрим жадный алгоритм: пока есть рёбра, выбираем вершину максимальной степени, добавляем её в покрытие $S$, удаляем эту вершину и все инцидентные ей рёбра. Доказать, что $|S|\le (1+\ln n)\,\mathrm{OPT}$ (натуральный логарифм, $n=|V|$), где $\mathrm{OPT}$ — размер минимального вершинного покрытия. Построить граф, на котором этот алгоритм даёт ответ, хуже оптимального в 100 раз.

  \paragraph{Алгоритм.}
  Повторять: взять вершину максимальной текущей степени, добавить в $S$, удалить её и все инцидентные рёбра. Остановиться, когда рёбер не осталось. Вернуть $S$.

  \paragraph{Корректность.}
  На каждом шаге удаляются все рёбра, инцидентные выбранной вершине. Алгоритм останавливается только когда рёбер не осталось, то есть каждое исходное ребро было удалено на одном из шагов. Значит, каждое ребро имеет хотя бы одну инцидентную вершину в $S$, следовательно $S$ — вершинное покрытие.

  \paragraph{Гарантия $1+\ln n$.}
  Рассмотрим оптимальное покрытие размера $k$. На любом шаге среди оставшихся рёбер найдётся вершина, покрывающая не меньше доли $1/k$ этих рёбер (иначе $k$ вершин не покрыли бы все оставшиеся рёбра). Жадный выбор берёт вершину с не меньшим числом инцидентных оставшихся рёбер, значит, за шаг удаляется как минимум доля $1/k$ оставшихся рёбер. После $t$ шагов число рёбер не превосходит $m_0\,(1-1/k)^t\le m_0\,e^{-t/k}$, где $m_0\le n(n-1)/2$. Отсюда достаточно $t\le k\,(1+\ln n)$, то есть $|S|\le (1+\ln n)\,k$.

  \paragraph{Пример с большим отрывом.}
  Есть семейство двудольных графов, на которых описанный жадный алгоритм даёт отношение порядка $\ln n$ к оптимальному. Идея: разбить правую долю на группы, убывающие геометрически, и соединять их так, чтобы вершины с наибольшей текущей степенью на каждом шаге были плохим выбором, тогда алгоритм делает примерно $\ln n$ шагов, а оптимум мал. Для достаточно больших $n$ это даёт пример, где жадный ответ хуже оптимального как минимум в 100 раз.

  \section*{Задача 0.6. Локальные улучшения для Max-Cut — остановка и гарантия 1/2}

  \paragraph{Условие.}
  Дано: неориентированный граф $G=(V,E)$. Алгоритм: начать с произвольного разбиения $V=V_1\cup V_2$, затем пока найдётся вершина $v$, у которой соседей в её текущей доле больше, чем соседей в другой доле, переместить $v$ в другую долю. Вернуть разбиение, когда таких вершин не останется. Доказать, что алгоритм останавливается и полученный разрез содержит не меньше половины всех рёбер.

  \paragraph{Корректность (остановка).}
  При каждом переносе вершины $v$ число рёбер разреза строго увеличивается: рёбра от $v$ к её текущей доле становятся пересекающими разрез и их больше, чем рёбер к противоположной доле, которые перестают пересекать разрез. Поскольку число рёбер конечно и размер разреза растёт на целое число, процесс не может продолжаться бесконечно — алгоритм останавливается.

  \paragraph{Гарантия 1/2.}
  В состоянии покоя у каждой вершины $v$ число соседей по её доле не больше числа соседей по другой доле (иначе правило позволило бы переместить $v$). Значит, для каждой вершины хотя бы половина инцидентных ей рёбер пересекает разрез. Суммируя по всем вершинам, получаем, что суммарно рёбер, пересекающих разрез, считая с кратностью по концам, не меньше $\sum_v \deg(v)/2 = 2|E|/2=|E|$. Каждое ребро разреза посчитано дважды (по обоим концам), поэтому размер разреза не меньше $|E|/2$. Отсюда найденное решение содержит не менее половины всех рёбер графа.


\end{document}


